\section{Calibration and local roads}\label{sec:roads}

Assume that the exponential time-constant ball has a nontrivial flat
face.  We associate a nonnegative defect with every path, then minimize
the expected defect of roads that track either of two directions in
the cone of that face.  Two estimates begin the argument: there exist
local roads with sublinear expected defect at large scales, and the
infimum of the common period bound over admissible pairs is bounded
away from zero.  Both estimates will remain available when the road
rules are changed later.

\subsection{A dominated stationary calibration}
\label{sec:fb-calibration}

Let $h$ be the supporting functional equal to one on the hypothetical
face.  Write $\mathcal G_n=n^{-1}\Z^2$; rescaling the lattice
does not rescale its exponential edge weights.  For $z\in\mathcal G_n$,
the shift $\theta_z$ translates microscopic inputs by $nz$.

The stationary Busemann-measure construction of Damron and Hanson
\cite[Section~3, Propositions~3.2 and~3.4, and Theorem~3.5]{DH2014}
provides, on a stationary extension of the weight space, an additive
cocycle $\mathcal B(x,y)$ with
\[
 |\mathcal B(x,y)|\le T(x,y),\qquad
 \E \mathcal B(0,x)=h(x),\qquad x,y\in\Z^2.
\]
The construction allows the prescribed supporting functional $h$ and
requires no differentiability of the norm.  Its moment and
subcritical-zero-mass assumptions hold for exponential weights.
For mean identification, Damron and Hanson use Hoffman's argument
\cite{Hoffman2008}, developed from an averaging argument of Garet and
Marchand \cite{GaretMarchand2005}, with a presentation inspired by
Gou\'er\'e's Lemma~2.6 \cite{Gouere2007}.

Condition $\mathcal B$ on the full edge environment.  This preserves
additivity, domination and the mean.  The resulting cocycle also has a
translation-covariant version: stationarity gives the conditional
expectation identity for each fixed translation, and countability of
the vertices and translations gives a common exceptional null set.
After rescaling, we therefore have an edge-measurable cocycle $B_n$ on
$\mathcal G_n$ satisfying
\begin{equation}
 \begin{split}
 B_n(x,z)&=B_n(x,y)+B_n(y,z),\\
 |B_n(x,y)|&\le T(nx,ny),\\
 \E B_n(x,y)&=n h(y-x),\\
 B_{n,\omega}(x+z,y+z)&=B_{n,\theta_z\omega}(x,y).
 \end{split}
 \label{eq:fb-calibration}
\end{equation}
These identities persist after adjoining independent private inputs.
For a path $\gamma$ from $x$ to $y$, put
\[
 P_n(x)=B_n(0,x)-nh(x),\qquad
 Q_n(\gamma)=\tau(\gamma)-nh(y-x)
\]
and define its \emph{defect} by
\begin{equation}
 D_n(\gamma)
 =Q_n(\gamma)-P_n(y)+P_n(x)
 =\tau(\gamma)-B_n(x,y)\ge0.
 \label{eq:fb-defect}
\end{equation}
Defect is additive under concatenation and nonnegative on every
forward subpath, with repeated edge occurrences counted with
multiplicity.  This pathwise accounting will allow endpoints to be
chosen after the environment has been inspected.

For cuts inside an edge, interpolate $B_n(0,\cdot)$ affinely and charge
the traversed fraction of the weight.  Each orientation of the edge
has nonnegative defect, so \eqref{eq:fb-defect} also holds for these
partial traversals.  Passage times between lattice vertices are
unchanged.

\subsection{The class of local randomized roads}
\label{sec:fb-roads}

Choose rational $v,w$ in the cone of the face, close to its two
endpoints, and a rational $0<s<1/4$.  Set
\[
 A=v+sw,\qquad B=w+sv,\qquad r=s/1000.
\]
We restrict to $n$ for which $nA,nB$ are lattice vectors.  All road
tracking distances and input-dependence radii are measured in the fixed
oblique norm
\[
 \|x\|_* =\bigl|[v\ w]^{-1}x\bigr|_\infty.
\]
In these coordinates $A=(1,s)$ and $B=(s,1)$.  In particular,
\begin{equation}
 \bigl|[A\ B]^{-1}x\bigr|_\infty\le2\|x\|_* ,
 \qquad
 \|x\|_*\ge\tfrac34\bigl|[A\ B]^{-1}x\bigr|_\infty.
 \label{eq:fb-coordinate-comparison}
\end{equation}
The conversion to $(v,B)$ coordinates has operator norm at most $2$
as well.  Thus changes between these oblique coordinates have uniform
constants.  Conversion to microscopic lattice coordinates may depend
on the fixed face: restrict the neighborhoods of its endpoints once
and choose $K_{\mathrm{lat}},M\ge1$ such that, for the standard basis
vectors $e_1,e_2$,
\begin{equation}
 \max_{i=1,2}\|e_i\|_*\le K_{\mathrm{lat}},
 \qquad |A|_2,|B|_2\le M.
 \label{eq:fb-geometry-constants}
\end{equation}
The constants depend only on the fixed face.  Such uniform choices
are possible because its endpoints, lying on $h=1$, are linearly
independent.

\begin{definition}[Admissible road]
\label{def:fb-road}
For $d\in\{A,B\}$, an admissible $d$-road is a measurable oriented
two-sided lattice walk together with a continuous nondecreasing nominal
parameter on its edge occurrences.  The nominal parameter tends to
$-\infty$ and $+\infty$ at the two ends.  Every compact parameter
interval contains finitely many edge occurrences.  A point with nominal
parameter $t$ satisfies $\|x-td\|_*\le r$.

The construction uses the edge environment and a private product field
$U=(U_\iota)_{\iota\in\mathcal I}$ independent of it. Each coordinate
has a spatial location and, when needed, a road-position or auxiliary
type label. For every $z\in\mathcal G_n$, the shift $\theta_z$ acts
measurably by translating the spatial and road-position indices together
and preserves the product law. Thus $(\theta_z\omega,\theta_zU)$ has
the same law as $(\omega,U)$. Marginals may depend on a type left
unchanged by translations. The road's restriction to
a parameter interval $I$, including all cuts and occurrences, is a
function of inputs at $\|\cdot\|_*$-distance at most $10$ from
$\{td:t\in I\}$.  The rule commutes with translation by $d$ and a unit
shift of nominal parameter.  At an integer parameter level use the
first occurrence of that level as the marked cut $X_i$.  In particular,
\begin{equation}
 X_i(\omega,U)
 =id+X_0(\theta_{id}\omega,\theta_{id}U).
 \label{eq:fb-cut-covariance}
\end{equation}
The cost between consecutive marked cuts has finite expectation.
Translated copies are obtained by translating this same rule, using
fresh private inputs when independence of the private randomizations
is required.
\end{definition}

A pause at an integer level belongs to the following period under
this convention; any other covariant marking gives the same accounting.
Locality applies to private inputs as well as edge weights.  In
particular, it does not permit a global random phase shared by distant
parts of a road.

Measurability of a parametrized walk is understood in the following
standard Borel coding. In a compact parameter window, record the finite
ordered list of labeled closed subintervals of oriented edge occurrences,
their endpoint fractions in the original edges, and their continuous
nominal maps after linear reparametrization to $[0,1]$. Use the usual
Borel structure of $C([0,1])$ with its uniform topology, together with
the discrete edge and occurrence labels. Marked occurrences are part
of this local record, and repeated visits retain distinct labels. This
coding also applies to every bounded road portion below.

\begin{lemma}[Random-cut period identity]
\label{lem:fb-random-cuts}
For any admissible $d$-road,
\begin{equation}
 \E [P_n(X_{i+1})-P_n(X_i)]=0,
 \qquad
 \E D_n(\gamma[i,i+1])
 =\E Q_n(\gamma[i,i+1])\ge0.
 \label{eq:fb-period-identity}
\end{equation}
Here $\gamma[i,i+1]$ denotes the oriented road portion between its
marked cuts.
\end{lemma}

\begin{proof}
By the joint stationarity of the edge and private inputs, it is enough
to consider $i=0$.  Combining the
centered cocycle identity with \eqref{eq:fb-cut-covariance} gives
\[
 P_{n,\omega}(X_1)
 =P_{n,\omega}(d)
  +P_{n,\theta_d\omega}
       \bigl(X_0(\theta_d\omega,\theta_dU)\bigr).
\]
The first term has mean zero; the second has the law of $P_n(X_0)$.
All terms are integrable because each cut lies on one of finitely many
deterministic edges, where domination bounds its calibration value by
a finite sum of weights and a deterministic linear term.  Taking
expectations proves the first identity, and \eqref{eq:fb-defect}
proves the second.  The mean of $P_n(X_0)$ itself need not vanish.
\end{proof}

For a pair of admissible roads define its common period bound as
\[
 \beta=\max\bigl\{
 \E Q_n(\gamma_A[0,1]),
 \E Q_n(\gamma_B[0,1])\bigr\},
\]
and let $b(n)$ be the infimum over admissible pairs.  The period
identity gives $b(n)\ge0$.  More generally, a forward portion between
levels $u\le v$ has defect bounded by the sum of the complete-period
defects whose parameter intervals meet $[u,v]$.  Nonnegativity makes
this a pathwise bound, valid for environment-selected endpoints.  We
will repeatedly use it to bound a random road portion by a fixed
number of periods.

\subsection{Initialization and a strictly positive infimum}
\label{sec:fb-bounds}

The two estimates for $b(n)$ use the same elementary fact: a long
crossing cannot be very cheap, uniformly over all choices of the
crossing path.  It first localizes geodesics between nearby anchors;
later it guarantees a saving when all weights in a small region are
decreased.

\begin{lemma}[No long cheap paths]
\label{lem:fb-cheap-paths}
Let $V$ be a finite set of vertices of $\Z^2$.  For every integer
$R\ge1$ and $t\ge0$,
\begin{equation}
 \PP \left(
 \begin{array}{c}
 \text{some self-avoiding path starting in $V$ has}\\[-2pt]
 \text{at least $R$ edges and cost at most $t$}
 \end{array}\right)
 \le \frac83 |V|e^{5t}2^{-R}.
 \label{eq:fb-cheap-bound}
\end{equation}
Consequently there are numerical constants $c,C>0$ such that, for
$z\ne0$ and $R\ge40|z|_1$, the probability that a geodesic from $0$ to
$z$ has at least $R$ edges is at most $Ce^{-cR}$.
\end{lemma}

\begin{proof}
A fixed vertex starts at most $4\cdot3^{m-1}$ self-avoiding paths
of length $m$.  For each path, independence and the identity
$\E e^{-5\tau_e}=1/6$ give
\[
 \PP (\tau(\gamma)\le t)\le e^{5t}6^{-m}.
\]
Sum over paths and over $m\ge R$ to obtain
\eqref{eq:fb-cheap-bound}.

For the geodesic estimate, compare with a deterministic coordinate
path of length $\ell=|z|_1$ and cost $S$.  Its exponential moment gives
$\PP (S>t)\le2^\ell e^{-t/2}$.  Since weights are positive, a
geodesic is self-avoiding.  If it has at least $R$ edges and $S\le t$,
it is counted by \eqref{eq:fb-cheap-bound} with $V=\{0\}$.  Taking
$t=(\log2)R/10$ makes both probabilities exponentially small in $R$
when $R\ge40\ell$.  The same estimate excludes arbitrarily long
paths of bounded cost.  Comparing with $S$ therefore also proves
existence of a finite minimizing path.
\end{proof}

\begin{proposition}[Sublinear initialization]
\label{prop:fb-upper}
For fixed $s,v,w$ there are admissible pairs of roads with
\begin{equation}
 b(n)\le C_s\sqrt n\log(2+n)\le C'_s n^{3/4}
 \label{eq:fb-upper}
\end{equation}
for all sufficiently large multiples of a fixed integer.  The constants
may depend on $s$; the exponent $3/4$ does not.
\end{proposition}

\begin{proof}
We concatenate geodesics between closely spaced deterministic anchors,
restricting each geodesic to a small ball to make the rule local.
Choose an integer $m$ large enough that, with $\delta=1/m$,
\[
 \delta\max(|A|_1,|B|_1)\le r/(1000K_{\mathrm{lat}}).
\]
Increase the fixed divisibility requirement so that $n\delta A$ and
$n\delta B$ are lattice vectors.  For $d=A$ or $B$, join successive
deterministic anchors $k\delta d$ by a minimizing path restricted to
the lattice portion of the $\|\cdot\|_*$-ball of scaled radius $r/4$
about the first anchor.  Each microscopic edge has $\|\cdot\|_*$-length
at most $K_{\mathrm{lat}}/n$, so a coordinate path between the anchors stays within
$K_{\mathrm{lat}}\delta|d|_1\le r/1000$ of the first anchor.  The restricted minimum
therefore exists.  Use a translation-covariant tie rule.

Restriction changes the minimum only if the unrestricted geodesic
leaves the ball.  Such a geodesic has at least
$rn/(4K_{\mathrm{lat}})$ edges, so Lemma~\ref{lem:fb-cheap-paths},
with integer rounding, bounds this event by $Ce^{-c_s n}$ for large
$n$.  On that event the increase is at most the coordinate-path cost
$S_k$.  Since $(\E S_k^2)^{1/2}\le C_s n$, Cauchy--Schwarz
bounds the expected increase by $C_s n e^{-c_s n/2}$.

For each patch, Lemma~\ref{model:metric} bounds the mean unrestricted
cost above its norm value by $C_s\sqrt n\log(2+n)$.  The latter value
is $\mu(n\delta d)=n\delta h(d)$ because $d$ lies in the cone of the
face.  Summing over the $m$ patches in a unit period proves the first
inequality in \eqref{eq:fb-upper}; increasing $C'_s$ gives the second.

Parametrize each patch monotonically from $k\delta$ to
$(k+1)\delta$, for example proportionally to its number of edge
occurrences.  Its $\|\cdot\|_*$-distance from its nominal location is
at most $r/4+\delta\|d\|_*\le r/4+K_{\mathrm{lat}}\delta|d|_1<r$.
The inputs for each patch lie within this same distance of every
nominal point in the patch.  Thus the radius-ten locality condition
holds.  The rule commutes with translation by $d$, and each period
has finite expected cost.  This proves admissibility.
\end{proof}

\begin{proposition}[Uniform positive lower bound]
\label{prop:fb-positive}
There is $c_0>0$, depending only on the fixed flat and the dependence
radius $10$, such that
\begin{equation}
 b(n)\ge c_0
 \label{eq:fb-positive}
\end{equation}
for all sufficiently large admissible $n$.  The constant $c_0$ is
independent of the road rules, of sufficiently small $s$, and of the
nearby rational directions.  The lower threshold on $n$ may depend
on $s$ and the divisibility requirements.
\end{proposition}

\begin{proof}
Two routes around a large square have small expected defects if
$\beta$ is small.  We decrease weights near the middle of one route,
leaving the other route and their common endpoints fixed.  The first
route then becomes a cheaper competitor, while the change in law is
small enough to transfer that saving back to the original environment.

\emph{The two routes.}
Fix an admissible pair with common bound $\beta$.  Write
$(\xi,\eta)=[A\ B]^{-1}x$ for its oblique coordinates, and put
$\rho=2r$.  Choose the fixed even integer
\[
 L=200.
\]
Use four translated roads, with nominal lines
\[
 \begin{aligned}
 \eta=0,\quad\eta=L
 &\quad\text{for the two $A$-roads},\\
 \xi=0,\quad\xi=L
 &\quad\text{for the two $B$-roads}.
 \end{aligned}
\]
They use the same physical edge weights and independent private
inputs.  Their strips have oblique half-width at most $\rho$.

At each nominal corner, determine the intersection from parameter
windows of half-width $4\rho$ on the two incident roads.  One
road crosses the square of coordinate half-width $2\rho$ from left
to right inside its horizontal strip, and the other crosses from
bottom to top inside its vertical strip.  The elementary planar
crossing property for a rectangle gives an intersection.  Nearest-
neighbor lattice edges can meet only at a common vertex or along a
common edge.  If $K_{\mathrm{lat}}/n\le r$, a common vertex can be chosen with
coordinate distance at most $2\rho$ from the corner and with both
nominal parameters within $3\rho$ of their corner values: moving
from an interior intersection to an endpoint of its edge costs at
most $2K_{\mathrm{lat}}/n\le\rho$ in coordinates, after which tracking costs at
most another $\rho$ in nominal parameter.  The chosen vertex remains
in the inspected windows.  Select the lexicographically first eligible vertex
and then the first pair of eligible occurrences.  These finite choices
give a measurable local corner rule.

The bottom and right sides, traversed forward between these selected
corners, form $R_1$; the left and top sides form $R_2$.  Their initial
and terminal vertices agree.  Since $3\rho<1$, every side used is a
subpath of the periods with indices $-1,0,\ldots,L$.  Nonnegative
defects and Lemma \ref{lem:fb-random-cuts} give, with
$C_R=2(L+2)$,
\begin{equation}
 \E D_n(R_j)\le C_R\beta,\qquad j=1,2,
 \qquad
 |\tau(R_1)-\tau(R_2)|\le D_n(R_1)+D_n(R_2).
 \label{eq:fb-route-defects}
\end{equation}

\emph{A region traversed by only one route.}
Consider the deterministic middle-side parallelogram
\[
 S=\{(L/2)A+uA+vB:|u|\le2,\ |v|\le2\},
\]
and let $\mathcal E_n(S)$ be the edges having both endpoints in
$S\cap\mathcal G_n$.  The entire input support of $R_2$, including
its selected endpoints and the middle corner, is outside $S$.
Indeed that support is contained in the $\|\cdot\|_*$-neighborhood of radius $10$
of the deterministic union $\mathcal C$ of the left and top nominal
segments and all incident nominal corner windows of parameter
half-width $4\rho$.  The coordinate distance from $S$ to $\mathcal C$ is at least
$L/2-2-4\rho$.  By \eqref{eq:fb-coordinate-comparison}, its
$\|\cdot\|_*$-distance is at least
$\tfrac34(L/2-2-4\rho)>10$.  The same separation holds for the unused corner window.
Thus changing only weights in $\mathcal E_n(S)$
leaves $R_2$ and its endpoints unchanged, even when its rule and
corner selections are recomputed.

The bottom-road portion with nominal parameters between
$L/2-1$ and $L/2+1$ belongs to $R_1$.  For sufficiently large $n$,
every edge meeting this portion has both endpoints in $S$.  Discarding
the first and last partial edges leaves endpoints with coordinate
separation at least
$2-2\rho-4K_{\mathrm{lat}}/n\ge1$.  Each microscopic edge has
$(A,B)$-coordinate length at most $2K_{\mathrm{lat}}/n$, so the resulting
lattice path has at least $n/(2K_{\mathrm{lat}})$ edges, up to integer
rounding.  Loop erasure preserves the endpoints and their separation,
so the resulting self-avoiding path obeys the same length bound and
has no greater cost.

There are at most $C_V n^2$ vertices in $S\cap\mathcal G_n$, where
one may take $C_V=(8M+3)^2$.  Apply Lemma
\ref{lem:fb-cheap-paths} with
\[
 R=\lfloor n/(2K_{\mathrm{lat}})\rfloor,\qquad
 a_*=(\log2)/(40K_{\mathrm{lat}}),\qquad t=a_*n.
\]
For $n\ge4K_{\mathrm{lat}}$ the resulting bound is
\begin{equation}
 \PP \left(
 \text{the cost of $R_1$ charged to $\mathcal E_n(S)$ is below $a_*n$}
 \right)
 \le q_n:=C n^2e^{-cn},
 \label{eq:fb-middle-saving}
\end{equation}
The constants $c,C>0$ depend only on $K_{\mathrm{lat}},M$.  Uniformity
in the road rule follows because the cheap-path event includes every
self-avoiding path starting at every vertex of $S$.

\emph{Perturbation and transfer of the saving.}
Keep the private inputs fixed.  For $0<\varepsilon\le1/2$, put
$a_n=1-\varepsilon/n$ and multiply every weight in
$\mathcal E_n(S)$ by $a_n$.  Let $\PP $ be the original joint
law of weights and private inputs, and $\mathbb Q$ its pushforward
under this map.  Write $m_n=|\mathcal E_n(S)|$.  Since each vertex
is charged at most two undirected edges,
$m_n\le C_E n^2$, where $C_E=2C_V$.  The relative entropy of a scaled
exponential variable with respect to an unscaled one is
$a_n-1-\log a_n$.  Product additivity, followed by Pinsker's
inequality \cite[Theorem~31, order one]{vanErvenHarremoes2014}, with total
variation taken as one half of the $L^1$ distance, gives
\begin{equation}
 \begin{split}
 D_{\mathrm{KL}}(\mathbb Q\Vert\PP )
 &=m_n(a_n-1-\log a_n)\le C_E\varepsilon^2,\\
 \|\mathbb Q-\PP \|_{\mathrm{TV}}
 &\le C_{\mathrm{TV}}\varepsilon,
 \qquad C_{\mathrm{TV}}=\sqrt{C_E/2}.
 \end{split}
 \label{eq:fb-total-variation}
\end{equation}
Here we used $-x-\log(1-x)\le x^2$ for $0\le x\le1/2$.

Set $\lambda=a_*\varepsilon/2$.  Let $E$ be the event that the cost
of the current $R_2$ exceeds passage time between its current endpoints
by at least $\lambda$.  This event uses only the current $R_2$, its
endpoints and the passage metric.  In the perturbation coupling,
outside the event in \eqref{eq:fb-middle-saving}, the old $R_1$ saves at least
$a_*\varepsilon$.  It still joins the endpoints of the unchanged
$R_2$.  If its old cost exceeded that of $R_2$ by at most $\lambda$,
it is therefore a competitor cheaper by at least $\lambda$ after
perturbation.  Equation~\eqref{eq:fb-route-defects} and Markov's
inequality give
\[
 \mathbb Q(E)
 \ge1-q_n-
 \PP \bigl(D_n(R_1)+D_n(R_2)>\lambda\bigr)
 \ge1-q_n-\frac{2C_R\beta}{\lambda}.
\]
Only the old $R_1$ is used as a competitor; its recomputed rule may
produce a different path.  Transferring $E$ to the original law by
\eqref{eq:fb-total-variation} yields
\[
 \PP (E)
 \ge1-q_n-C_{\mathrm{TV}}\varepsilon
       -\frac{2C_R\beta}{\lambda}.
\]
In the original environment domination implies
$D_n(R_2)\ge\lambda$ on $E$, also for its random endpoints.
Hence
\[
 C_R\beta\ge\lambda\PP (E)
 \ge\lambda(1-q_n-C_{\mathrm{TV}}\varepsilon)-2C_R\beta.
\]
Fix $\varepsilon>0$ with $C_{\mathrm{TV}}\varepsilon\le1/8$ and
then take $n$ large enough that $q_n\le1/8$.  Rearranging yields
\[
 \beta\ge \frac{a_*\varepsilon}{8C_R}=:c_0>0.
\]
The constants in $c_0$ depend only on the uniform geometry bounds
and the fixed locality radius.  Taking the infimum over admissible
pairs proves \eqref{eq:fb-positive}.  The transfer used the event $E$,
so no comparison of calibrations in the two environments was needed.
\end{proof}

\section{Local gauges for nearly calibrated roads}
\label{found:sec-gauges}

The defect estimates above involve the generally nonlocal calibration
$P_n$.  The next step replaces its increments by increments of a
\emph{local gauge}: an endpoint potential computed from nearby inputs.
For endpoints with specified road labels, the corrected cost is
$Q_n(\gamma)-g(y)+g(x)$.  These corrections telescope whenever
successive pieces use the same endpoint label.

The construction has three stages.  Short signed trips along crossing
roads estimate calibration increments.  A finite-cell projection turns
these local estimates into local potentials.  Interpolation then
controls all attachment occurrences in a bounded window at once.
Throughout, we estimate differences of calibration errors; we do not
need a local approximation to the absolute value of $P_n$.

Expectations include the roads' private randomness.  We may condition
throughout on a fixed finite collection of phases, provided every
hypothesis, including the period defect bound, holds uniformly over
the conditioned phase types.
We retain the norm $\|\cdot\|_*$ and geometry constants of
Section~\ref{sec:fb-roads}. Conversion to $(A,B)$ or $(v,B)$ coordinates
has operator norm at most $2$ uniformly for $0<s<1/4$; the scaled
edge length is at most $K_{\rm lat}/n$, where $K_{\rm lat}$ may depend
on the fixed flat. Suppress the subscript $n$ in this section.
We write $Q(\gamma)=\tau(\gamma)-nh(y-x)$ for a path from $x$ to $y$,
$P(x)=B_n(0,x)-nh(x)$, and
\[
 D(\gamma)=Q(\gamma)-P(y)+P(x)\geq0.
\]
A signed trip is a concatenation of forward road intervals with
specified signs: both $Q$ and the increments of $P$ are summed with
those signs.  It estimates a potential difference.  Physical path
costs, by contrast, always charge every traversal positively.

\subsection{A bounded network of signed comparisons}

\begin{lemma}[Local increment comparisons]
\label{found:lem-network}
Let $A=v+sw$, $B=w+sv$ with $0<s<1/4$, and suppose two road rules have
tracking radius $r$ and dependence radius $10$ in $\|\cdot\|_*$, and
expected defect at most $\beta$ on each marked unit period.  There
exist numerical constants $r_0,H,C>0$ with the following property
whenever $r<r_0$ and $n\geq K_{\rm lat}/r$.

One can choose a random anchor $X_z$ near each point $z_1A+z_2B$,
$z\in\Z^2$, and, for $i=1,2$, a random variable $E_{z,i}$ such
that
\begin{equation}
 \E \left|E_{z,i}-P(X_{z+e_i})+P(X_z)\right|
 \leq C\beta.                                      \label{found:eq-increment}
\end{equation}
The variable $E_{z,i}$ uses only input cells at index distance at most
$H$ from $z$.  Here an input cell contains all microscopic weights and
private inputs assigned to the corresponding unit parallelogram.

One may also prescribe any collection of additional marked anchors
$X_{z,\lambda}$, each on a translated $A$- or $B$-road within a fixed
distance of $z_1A+z_2B$.  There are comparisons of the same kind between
any two such anchors at the same or neighboring indices.  Their constants
depend on that distance, but not on the number of labels $\lambda$.
\end{lemma}

\begin{proof}
We build a grid from the two road families, then attach any additional
anchor to that grid by a bounded signed trip.  Put an $A$-road on each
line $jB+\R A$ and a $B$-road on each
line $iA+\R B$.  Write physical coordinates in the basis $(A,B)$.
The coordinate-conversion bound implies that the first road lies in
$|x_2-j|\leq\delta$ and the second in $|x_1-i|\leq\delta$, where
$\delta=2r$.  Take $r_0=1/100$, so $8\delta<1/2$.

The portions whose nominal parameters are within $4\delta$ of the
nominal crossing connect opposite sides of a small rectangle.
The planar crossing property gives an intersection of the two portions.
An intersection lies in both strips and hence within coordinate
distance $\delta$ of the nominal crossing.  Nearest-neighbor paths
meet at a common vertex or along a common edge.  Each scaled edge has
$(A,B)$ coordinate length at most $2K_{\rm lat}/n\leq\delta$; a vertex
of that edge therefore has coordinate distance at most $2\delta$
from the crossing and nominal parameter distance at most $3\delta$
on either road.  It lies within the inspected portions.  Choose the
first such vertex occurrence along the
$A$ portion, breaking any remaining tie by its occurrence on the $B$
portion.  This choice uses only those two bounded portions, is covariant
under lattice translations, and specifies $X_{(i,j)}$ together with its
occurrence label on each road.  Loops and repeated visits cause no
ambiguity because occurrences, rather than only vertex names, are used.

The disjoint parameter windows order neighboring intersections along
the relevant road.  Define $E_{z,1}$ as the $Q$-cost of the intervening
forward $A$ portion, and define $E_{z,2}$ using the $B$ portion.
Each portion lies in a fixed number of marked periods, so additivity
and nonnegativity bound its expected defect by $C\beta$.  This proves
\eqref{found:eq-increment}.

Tracking and locality put all inputs in a fixed number of unit
parallelograms.  Assign boundary weights by a fixed half-open convention
and private inputs to the cell of their nominal location.  Since each
microscopic input is assigned to one cell, the cell variables are
independent.

To attach an additional anchor, follow its road with either sign to
a nearby transverse network road, then follow that road to a network
intersection.  A bounded number of network intervals reaches any
other anchor under consideration.  All nominal distances are bounded,
and crossings use the same local selection rule.  The error on each
signed interval is its signed defect, whose absolute value is bounded
by complete-period defects.  The triangle inequality gives the stated
comparison.  Its constant depends on the allowed anchor distance,
independently of the number of other labels.
\end{proof}

\subsection{An explicit finite-cell lemma}

We now isolate the step from local increment estimates to a local
potential.  Only independence of coarse cells is needed.  A cell may
contain arbitrarily many microscopic variables, so the constants below
do not depend on lattice scale or on the fineness of a fixed row mesh.

\begin{lemma}[Local increments admit a local potential]
\label{found:lem-finite-cell}
Let $(\eta_z)_{z\in\Z^2}$ be independent random elements and let
$V_{z,\lambda}\in L^1$, where $\lambda$ belongs to an arbitrary label
set with a distinguished label $0$.  Set
$\mathcal F_I=\sigma(\eta_z:z\in I)$ and
$Q_H(z)=z+[-H,H]^2\cap\Z^2$, with integer $H\geq2$.
Suppose that whenever $|z-z'|_1\leq1$, and for every pair of labels
in question, there is an $\mathcal F_{Q_H(z)}$-measurable variable
$E_{(z,\lambda),(z',\lambda')}$ such that
\begin{equation}
 \left\|E_{(z,\lambda),(z',\lambda')}
       -V_{z',\lambda'}+V_{z,\lambda}\right\|_1
 \leq\varepsilon.                              \label{found:eq-abstract-hyp}
\end{equation}
Then there are random variables $g_{z,\lambda}$, measurable with respect
to $\mathcal F_{Q_{R+2H}(z)}$, where $R=8H+20$, such that
\begin{align}
 \E g_{z,\lambda}&=\E V_{z,\lambda},\label{found:eq-mean}\\
 \left\|g_{z',\lambda'}-g_{z,\lambda}
       -V_{z',\lambda'}+V_{z,\lambda}\right\|_1
 &\leq C_H\varepsilon.                         \label{found:eq-potential}
\end{align}
The constant $C_H$ depends only on $H$, not on the number of labels.
Translation covariance is preserved whenever the true increment field
$V_{z',\lambda'}-V_{z,\lambda}$ is covariant under measure-preserving input shifts
and the deterministic means have the required covariance.
\end{lemma}

\begin{proof}
We decompose a local increment according to the finite sets of cells
on which it depends.  Each component can be integrated locally; an
exterior detour controls the error where that integration is cut off.

For a nonempty finite $S\subset\Z^2$, define
\[
 \Pi_S F=\sum_{I\subset S}(-1)^{|S|-|I|}
                 \E [F\mid\mathcal F_I].
\]
These are finite-product ANOVA (Hoeffding) projections; see
Efron and Stein \cite[Section~2]{EfronStein1981} for the classical
square-integrable decomposition. Here we use only finite
inclusion--exclusion identities in $L^1$. Each projection has mean
zero and $L^1$ operator norm at most $2^{|S|}$.
If $F$ is measurable with respect to a finite cell set $N$, then
$\Pi_SF=0$ unless $S\subset N$, and
\begin{equation}
 F=\E F+\sum_{\varnothing\ne S\subset N}\Pi_SF.
                                                        \label{found:eq-anova}
\end{equation}
Cancellation and independence prove both identities directly in
$L^1$; no infinite or $L^2$ expansion is involved.

A comparison supported on $Q_H(z)$ has components only on sets of
at most $m=(2H+1)^2$ cells and diameter at most $2H$.  Accordingly,
let $\mathscr S$ be the collection of nonempty finite
sets $S$ with $|S|\leq m$ and $\operatorname{diam}_\infty S\leq2H$.
Let $k(S)$ be the lexicographically least point of $S$, put
\[
 U_S=Q_R(k(S)),\qquad
 z_S=k(S)+(R+4H+10)e_1,
\]
and use $z_S$ as a reference beyond the cutoff $U_S$.  Define
\begin{equation}
 g_{z,\lambda}=\E V_{z,\lambda}
    +\sum_{\substack{S\in\mathscr S\\z\in U_S}}
       \Pi_S\bigl(V_{z,\lambda}-V_{z_S,0}\bigr).
                                                        \label{found:eq-gauge-formula}
\end{equation}
There are only a bounded number $N_H$ of summands: $k(S)$ must lie in
$Q_R(z)$ and $S\subset Q_{2H}(k(S))$.  For example
$N_H\leq(2R+1)^2 2^{(4H+1)^2}$ suffices.  Every summand is measurable
with respect to cells in $Q_{R+2H}(z)$ and has mean zero.  This proves
locality and \eqref{found:eq-mean}.

\emph{The cutoff error.}
Suppose $z\in U_S$ but a neighbor $z'$ lies outside.  There is a
lattice path from $z$ to
$z_S$ of length at most $20(R+4H+10)$ lying outside
$Q_{R-2}(k(S))$: follow the boundary of $Q_R(k(S))$ to its eastern
side and then move east.  Every local support $Q_H(u)$ along this path
is disjoint from $S\subset Q_{2H}(k(S))$.  If necessary first switch
from label $\lambda$ to label $0$ at $z$.  By
\eqref{found:eq-abstract-hyp}, telescoping and the projection norm,
\begin{equation}
 \left\|\Pi_S(V_{z,\lambda}-V_{z_S,0})\right\|_1
 \leq [20(R+4H+10)+1]2^m\varepsilon.
                                                        \label{found:eq-boundary}
\end{equation}
Each local comparison has zero $\Pi_S$ projection because its support
misses $S$.  The connected exterior in two dimensions is used exactly
at this step.

\emph{Recovery of an increment.}
Consider an edge $(z,\lambda),(z',\lambda')$ from the hypothesis
and write $W=V_{z',\lambda'}-V_{z,\lambda}$ and $E$ for its local
comparison.  The centered variables satisfy
$\|W-\E W-E+\E E\|_1\leq2\varepsilon$.
Every nonempty subset of $Q_H(z)$ contributes to both endpoint
gauges, since $R>H+1$.  A set contributing at both endpoints gives
$\Pi_SW$, within $2^m\varepsilon$ of $\Pi_SE$; this latter projection
vanishes when the set is not contained in $Q_H(z)$.  A set contributing
at only one endpoint is controlled by \eqref{found:eq-boundary}.
Thus \eqref{found:eq-anova}, the centered error bound above and at
most $2N_H$ component errors prove \eqref{found:eq-potential}.
For a same-index label change there are no cutoff terms.

Finally, the choices $k(S)$, $U_S$ and $z_S$ commute with index
translations.  The random summands in \eqref{found:eq-gauge-formula}
contain only differences of $V$, so a random common additive constant
under translation cancels.  Covariance of the increments and the
deterministic means therefore passes to the gauge.
\end{proof}

\begin{remark}[Random-anchor means]
\label{found:rem-means}
In the application $V_{z,\lambda}=P(X_{z,\lambda})$, the deterministic
mean in \eqref{found:eq-gauge-formula} need not be zero.  It is,
however, compatible with translation.  If
$X_{z+k,\lambda}(\omega)=t_k+X_{z,\lambda}(\theta_{t_k}\omega)$,
the centered cocycle identity and $\E P(t_k)=0$ give
$\E P(X_{z+k,\lambda})=\E P(X_{z,\lambda})$.
Thus it is a label mean, compatible with spatial covariance.  A common
deterministic constant may be subtracted from every label.
\end{remark}

\subsection{Interpolation and uniform envelopes}

Apply Lemmas~\ref{found:lem-network} and
\ref{found:lem-finite-cell} with $\varepsilon=C\beta$.  To extend the
result from marked anchors to arbitrary attachments, let $g_i$ be the
gauge at the $i$th marked cut of a road and put
\[
 C_i=Q(i,i+1)-g_{i+1}+g_i.
\]
Additivity and nonnegativity of defect give
\begin{equation}
 \E |C_i|\leq C\beta.                 \label{found:eq-slope}
\end{equation}
For an occurrence between the marked cuts $i$ and $i+1$, with nominal
parameter $t\in[i,i+1]$, set
\begin{equation}
 g(t)=g_i+Q(i,t)-(t-i)C_i.                  \label{found:eq-interpolation}
\end{equation}
Here $Q(i,t)$ includes the occurrence label.  During a pause the formula
is applied separately at successive occurrences; as before, a pause
at an integer belongs to the following period.  For any two ordered
occurrences in the same marked period, cancellation gives
\begin{equation}
 Q(t,u)-g(u)+g(t)=(u-t)C_i.                 \label{found:eq-short}
\end{equation}
Reversing a signed trip changes the sign of this identity.  Splitting
at integer levels shows that a prescribed window of length $\ell\le1$
has absolute corrected cost at most $\ell$ times the sum of $|C_i|$
over its at most two periods.  In particular, taking a supremum over
attachment occurrences costs no factor for the number of lattice
vertices in the window.

Writing $P(t)$ for the calibration at the physical occurrence,
\begin{equation}
 \bigl[g(t)-P(t)\bigr]-\bigl[g_i-P(i)\bigr]
       =D(i,t)-(t-i)C_i.                  \label{found:eq-interpolation-error}
\end{equation}
The right-hand side has expected absolute supremum at most $C\beta$,
because $0\le D(i,t)\le D(i,i+1)$.  Together with the anchor
increment estimates, this gives the simultaneous bounds needed for
randomly selected attachments.

\paragraph{Local extrema with variable endpoints.}
All parameter windows used in a local optimization are closed and
bounded. A bounded polygonal domain denotes its closed intersection
with the metric graph, subdivided at the finitely many endpoints of
point or segment intersections with its boundary. These conventions
justify the extrema below even when
an attachment lies inside an edge.

For fixed local inputs, a window contains finitely many labeled edge
occurrences. Positive edge weights allow loops in a physical competitor
to be deleted. Deletion preserves the later band and drawdown constraints,
because it retains a subsequence of the traversal order. There are
therefore finitely many combinatorial route types. For one type, let
$(a,b)\in[0,1]^2$ be the fractions specifying its two endpoints.
Membership in the closed parameter windows is a closed condition by
continuity of the nominal maps. Membership in the subdivided domain,
and the later band and drawdown conditions, are also closed: on each
linear segment the coordinate extrema occur at its endpoints, leaving
finitely many inequalities. Thus the feasible fractions form a compact
set. Partial passage costs and the interpolated gauges are continuous
in these fractions, including during a nominal pause. Each corrected
objective consequently attains its extrema on every nonempty type.

The extrema and a selected optimizer are measurable from the same local
inputs. To see this directly, enumerate the route types by a fixed order
of their relative edge and occurrence codes. For one type its constraints
can be written as $F(a,b)=0$, where $F\ge0$ is continuous in $(a,b)$
and measurable in the local record; write $c(a,b)$ for its continuous
cost. With the infimum of the empty set equal to $+\infty$, compactness
gives
\[
 \min_{F=0}c
 =\lim_{j\to\infty}\inf\{c(q):q\in\mathbb Q^2\cap[0,1]^2,
                                      \ F(q)<1/j\}.
\]
The right side is measurable. For an empty feasible set it is eventually
$+\infty$. When some type is feasible, take the first minimizing type,
then successively minimize the two endpoint coordinates on its compact
set of minimizers. The same formula makes this selector measurable, and
the relative order makes it commute with translations. If every type is
empty, return a fixed no-path symbol. The same argument applied to $-c$
handles suprema, and equality of two physical attachment points is one
more closed constraint for switch costs. A nonnegative envelope over an
empty family, including the negative part defining $U'$, is set to zero.

\begin{lemma}[Local envelope]
\label{found:lem-envelope}
Fix $m$ road labels and parameter windows of length at most $L$ contained
in a bounded macroscopic neighborhood of a deterministic cell.  There
is a nonnegative local random variable $U$ with
$\E U\leq C_{m,L}\beta$ which bounds all of the following:
absolute corrected signed costs within these windows; sums with a
bounded number of switches between these labels at intersections; and
the gauge differences for such switches at their common physical
vertex.  The support radius depends only on $L$, the
neighborhood size, and the original road locality radius.

If a physical path is confined to a prescribed deterministic finite
closed metric subgraph and has endpoints on two of the windows,
its corrected cost, including signed adjustments to prescribed center
anchors, is bounded below by $-U'$ for a local nonnegative variable
$U'$ with $\E U'\leq C_{m,L}\beta$.  Its support also includes
the prescribed subgraph.  This bound holds simultaneously for every such
path and every choice of its endpoints.
\end{lemma}

\begin{proof}
Define $U$ as a fixed multiple of the sum of the suprema of the listed
corrected costs and switch gauge differences, large enough for the
allowed number of switches.  It is local because the road windows
and their gauges are local.

For its expectation, compare with an auxiliary sum over the finitely
many relevant periods and anchors: include $|C_i|$, the defects
$D(i,i+1)$, and the absolute increment discrepancies in
\eqref{found:eq-potential}.  At a switch the two calibration values
refer to the same physical vertex and cancel.
Equations~\eqref{found:eq-short} and
\eqref{found:eq-interpolation-error} therefore bound $U$ by a constant
times this auxiliary sum.  Its expectation is $O(\beta)$, regardless
of how many microscopic occurrences lie in the windows.

For physical competitors, let $U'$ be the negative part of the minimum
endpoint-adjusted corrected cost in the prescribed subgraph.  This is
again a local measurable variable by the preceding endpoint-selection
argument. Domination gives $Q\ge\Delta P$ for
every competitor.  After endpoint corrections, only calibration-error
increments and signed corrected anchor adjustments remain in the
lower bound.  The same auxiliary sum controls their suprema in
expectation by $O(\beta)$.

The auxiliary estimates may be nonlocal; $U$ and $U'$ were defined
from the local costs themselves.  Thus the argument does not require
local pointwise dominators for the absolute errors $P-g$.
\end{proof}

\begin{lemma}[Linear-scale tail bound]
\label{found:lem-gauge-tail}
Suppose the period excess bound satisfies $\beta\leq C_s n^\alpha$
for some fixed $\alpha<1$.  Fix the neighborhood and finitely many
labels of Lemma~\ref{found:lem-envelope}.  After subtraction of one
common deterministic reference-label mean, a local random variable $V$ bounds the
absolute gauges at all those occurrences and the absolute raw signed
$Q$-costs of their bounded road adjustments.  For any fixed
$\alpha'\in(\max\{\alpha,1/2\},1)$, Lemma~\ref{model:metric} gives
\[
 \E V\leq C_s n^{\alpha'},\qquad
 \PP (V>\epsilon n)\leq C_s\epsilon^{-1}n^{\alpha'-1}.
\]
The constants may depend on the fixed neighborhood and label collection.
\end{lemma}

\begin{proof}
We bound the expectations of the local suprema appearing in the
statement, then take their sum as $V$.  The passage-time deficits
used to estimate these expectations need not be local.

Choose a fixed macroscopic box containing the road portions in question,
the bounded comparison trips used below and every lattice edge meeting
them.  Let $K_n^{\rm vert}$ be the
maximum of $(\mu(n(y-x))-T(nx,ny))_+$ over lattice vertices in that
box, and let $M_n$ be its maximum edge weight.  Put
\[
 C_\mu=2\max_{i=1,2}\mu(e_i),\qquad
 K_n=K_n^{\rm vert}+2M_n+C_\mu.
\]
The lower concentration estimate, summed over polynomially many
vertex pairs and then integrated, gives
$\E K_n^{\rm vert}\le C n^{\alpha'}$.  The exponential tail
and $O(n^2)$ edges give $\E M_n=O(\log(2+n))$.  Hence
$\E K_n\le Cn^{\alpha'}$.

The enlargement handles interior-edge endpoints as well.  Given a
road subinterval from $x$ to $y$, attach partial edges to endpoints
$x',y'$ of the incident lattice edges.  The resulting vertex-to-vertex
path costs at most the subinterval cost plus $2M_n$, while
\[
 \mu(n(y'-x'))\geq\mu(n(y-x))-C_\mu.
\]
The definition of $K_n^{\rm vert}$ therefore shows that every forward
road subinterval in the box satisfies $Q\geq-K_n$, because $\mu\geq h$.
For a complete period $q=Q(i,i+1)$ we therefore have
$\E |q|\leq\E q+2\E K_n
\leq\beta+2\E K_n$.
If $q=q_1+q_2+q_3$ is split at any two occurrences, then
$q_2\geq-K_n$ and $q_2\leq q+2K_n$.
Thus the supremum of absolute raw subinterval costs is bounded by
$|q|+2K_n$, including for intervals selected after inspecting the road.
The same conclusion holds for a bounded signed trip by addition.

For the gauges, each summand in \eqref{found:eq-gauge-formula} is a
projection of a potential difference.  Lemma~\ref{found:lem-network}
approximates that difference by a bounded signed trip with $L^1$
error $O(\beta)$.  The projection norm and the fixed number of terms
therefore give an $O_s(n^{\alpha'})$ bound.  Taking expectations of
the same comparison bounds the differences of deterministic label
means.  Interpolation by \eqref{found:eq-interpolation} extends the
bound to the windows.  After the common deterministic centering, the
sum of these local absolute suprema supplies $V$.  Markov's inequality
then gives the asserted tail.

For random phases, take the common centering constant to be the mean
of a specified fixed-phase reference anchor, and attach its translated
copy as an extra label near each site.  The signed-trip comparison
bounds the difference between any phase-conditioned anchor mean and
this reference by $C_s n^{\alpha'}$, uniformly in the phase.  The
centered cocycle identity makes the reference mean independent of the
site.  Thus centering remains deterministic and introduces no shared
random phase.
\end{proof}

\subsection{A common gauge in several exterior completions}

The preceding construction uses a calibration from one complete
environment.  We now show that its gauge on a preserved incoming road
can be approximated using only that road's own inputs.  This gives a
common endpoint correction when the exterior environment is completed
in several different ways.

\begin{lemma}[A gauge from the preserved incoming road]
\label{found:lem-common-gauge}
Fix deterministic tubes and a baseline road contained in one of them.
Write $\eta=(\eta_0,\eta_{\rm other})$ for all preserved tube weights,
where $\eta_0$ consists exactly of the weights in the baseline's own
tube and $\eta_{\rm other}$ contains the remaining preserved edges.
Let $\xi$ contain the baseline's private road seeds and a private
completion of the entire complement of its own tube.  Construct the
baseline once, using the actual weights in its own tube and these
private inputs everywhere else.  The resulting road is a function of
$(\eta_0,\xi)$ alone.

In each world $j$, retain all prescribed tube weights and use an iid
completion $\zeta^j$ elsewhere.  For every $j$, the private input
$\xi$ is independent of the joint tuple $(\eta,\zeta^j)$; different
worlds may share coordinates with one another.  Assume that every world
has the correct product weight law, uses the same measurable
calibration map, and supplies the local comparison network above with
period bound $\beta$.  Assume also that the shared baseline is
admissible in every world and has expected defect at most $\beta$ on
each marked period there.  The latter condition follows from the
random-cut period identity when the shared baseline rule has period
excess at most $\beta$.

Let $g_x^j$ be the gauge given by \eqref{found:eq-gauge-formula} at a
marked baseline anchor.  There is a local function $g_x^{\rm in}$ of
$(\eta_0,\xi)$ alone such that, in every world,
\begin{equation}
 \E |g_x^j-g_x^{\rm in}|\leq C\beta.
                                                        \label{found:eq-common-gauge}
\end{equation}
Its support radius and the constant $C$ are uniform in the number of
worlds.  After interpolation along any bounded baseline interval,
the same assertion holds with an expected supremum bound of order
$\beta$.  No identification of the calibrations in different worlds
is asserted.  In particular, incoming gauges for disjoint baseline
tubes with independent private completions and seeds are independent.
\end{lemma}

\begin{proof}
Use the same deterministic coarse cell partition in every world.
Within each cell, group the common tube coordinates, common private
coordinates, and additional coordinates of that world.  These cells
are independent.  Until we establish that a projection is common to
all worlds, $\Pi_S$ denotes the finite-cell projection in the product
space of the world under consideration.

We first replace each calibration difference in the gauge formula by
a common baseline cost.  Let $x$ be the anchor's coarse index, and
choose a marked anchor $y$ on the same baseline, a fixed integer
number $L_0$ of periods earlier.  Take $L_0$ large enough that $y$
lies outside every cutoff $U_S$ contributing at $x$ in
\eqref{found:eq-gauge-formula}, with an additional margin of $4H+20$.
The coordinate-conversion bound allows one choice of $L_0$ for all
$0<s<1/4$.

Let $Q_{y,x}$ be the forward baseline cost from $y$ to $x$.  The
traveled edges lie in the baseline's own preserved tube, and the path
and its marked occurrences depend only on $(\eta_0,\xi)$.  Thus
$Q_{y,x}$ is the same random variable in every world.  Adding the
period defects along this fixed number of periods gives
\[
 \|Q_{y,x}-P^j(X_x)+P^j(X_y)\|_1\leq C\beta.
\]

Fix a contributing set $S$.  Attach $y$ to the comparison network
locally near $y$, and connect that network anchor to the reference
anchor at $z_S$ by signed comparisons along the boundary of
$Q_R(k(S))$ and its exterior.  Such a
route can reach the square's boundary, follow it to the eastern side,
and then move east to $z_S$, just as in the cutoff argument leading
to \eqref{found:eq-boundary}.  The initial attachment is separated
from $S$ by the margin at $y$.  Along the remaining route, the local
supports $Q_H(u)$ miss $S\subset Q_{2H}(k(S))$, since
$R=8H+20$.  The route has bounded length: $L_0$ is fixed,
$k(S)\in Q_R(x)$, and $z_S-k(S)$ is fixed by the gauge construction.

Every local comparison on this route therefore has zero $\Pi_S$
projection.  Telescoping their calibration differences, using their
$O(\beta)$ errors and the bounded projection norm, yields
\[
 \|\Pi_S(P^j(X_y)-P^j(X_{z_S,0}))\|_1\leq C\beta.
\]
Combining this estimate with the preceding baseline estimate shows
that each summand of $g_x^j$ differs by $O(\beta)$ from
$\Pi_SQ_{y,x}$.

It remains to check that these projected baseline costs, including
the deterministic mean term, are common to all worlds.  The joint
marginal law of the world weights and the shared baseline is the same
in every world: the baseline uses its own tube weights together with
private inputs independent of that world's actual exterior.  Since
the calibration map is also the same, the mean
$m_x=\E P^j(X_x)$ is independent of $j$.  Define
\[
 g_x^{\rm in}=m_x+
       \sum_{\substack{S\in\mathscr S\\x\in U_S}}
                        \Pi_SQ_{y,x}.
\]
For any cell subset, conditioning $Q_{y,x}$ on its independent
$\eta_{\rm other}$ or $\zeta^j$ coordinates adds no information to
its $(\eta_0,\xi)$ coordinates.  Consequently, every conditional
expectation in $\Pi_SQ_{y,x}$ is a function only of that subset's
$(\eta_0,\xi)$ coordinates.  These conditional expectations, and
hence the projections themselves, can be chosen as the same
functions in every world.  They have no dependence on weights in
the other preserved tubes.

Only a fixed number of sets $S$ contributes, all within the fixed
support radius of \eqref{found:eq-gauge-formula}.  Summing the
projection errors proves \eqref{found:eq-common-gauge} and locality,
with constants independent of the number of worlds.  The displayed
construction also proves the independence assertion: disjoint
baseline tubes supply disjoint sets of product edge weights, their
private input collections are independent, and $m_x$ is deterministic.

Finally, interpolate the common anchor gauges using the common raw
baseline cost.  In \eqref{found:eq-interpolation}, the raw $Q$ terms
cancel when the two interpolations are subtracted.  For
$t\in[i,i+1]$, their discrepancy is therefore the linear
interpolation of the discrepancies at $i$ and $i+1$.  Its supremum
is bounded by the sum of those two anchor discrepancies.  A bounded
interval meets only a bounded number of periods, which proves the
stated expected supremum bound.
\end{proof}

\begin{remark}[Finite averaging of mesh offsets]
\label{found:rem-offsets}
Suppose the required row translations generate a finite group modulo
the coarse mesh periods.  Perform the construction for every
translated coarse mesh, and average the resulting gauge functions.
Each required translation permutes these finitely many functions,
so their average has exact covariance under the row translations.
All their supports lie in one fixed-radius ball.  By the triangle
inequality, averaging preserves every preceding $L^1$ bound without
a factor equal to the number of offsets.  Applying the incoming-road
construction separately for each offset and then averaging also
preserves Lemma~\ref{found:lem-common-gauge}.  The averaging is over
deterministic offsets and introduces no shared random mesh phase.
\end{remark}
